%! Polynomial Functions and Real Zeros — Unit 2, Lesson 2.3 — Guided Notes
\documentclass[10pt]{article}
\usepackage{precalculus-article}
\usepackage{precalculus-boxes}
\newcommand{\TallMath}[1]{$\displaystyle #1\rule[-1.4em]{0pt}{3.2em}$}

\begin{document}

\pageheader{Unit 2, Lesson 2.3}{Guided Notes}

\vspace{0.12in}

\begin{objectivebox}
By the end of this lesson, I will be able to\ldots
\begin{enumerate}[itemsep=2pt]
  \item Read the real \blank{1.6cm} of a polynomial from its factored form,
        and write a factored form from a list of zeros.
  \item Find the \blank{2.4cm} of each real zero and predict whether the graph
        \blank{1.7cm} the $x$-axis there or is \blank{1.9cm} to it.
  \item Find the \blank{1.6cm} of a polynomial from a table of values using
        successive differences.
\end{enumerate}
\end{objectivebox}

\vspace{0.1in}

\boxguard
\begin{vocabbox}
Fill in each term as we define it together.
\par\vspace{2pt}
\noindent\textbf{\textcolor{plum}{Zero of a function:}}\\[1pt]
\writeline\\[3pt]
\noindent\textbf{\textcolor{plum}{Root:}}\\[1pt]
\writeline\\[3pt]
\noindent\textbf{\textcolor{plum}{$x$-intercept:}}\\[1pt]
\writeline\\[3pt]
\noindent\textbf{\textcolor{plum}{Linear factor:}}\\[1pt]
\writeline\\[3pt]
\noindent\textbf{\textcolor{plum}{Multiplicity:}}\\[1pt]
\writeline\\[3pt]
\noindent\textbf{\textcolor{plum}{Tangent to the $x$-axis:}}\\[1pt]
\writeline
\end{vocabbox}

\vspace{0.1in}

\boxguard
\begin{hookbox}
Four questions, four different tests, one polynomial $p(x) = x^2 - x - 6$:

\smallskip
\textit{Solve $p(x) = 0$.} \quad \textit{Find the roots.} \quad
\textit{Find the zeros of $p$.} \quad \textit{Find the $x$-intercepts.}

\smallskip
What do all four questions have in common --- and what is the one answer?

\writelines{2}
\end{hookbox}

\vspace{0.1in}

\boxguard[30]
\begin{notesbox}{1. Four Names for One Idea}
A number $a$ is a \textbf{zero} of the function $p$ when
$p(a) = $ \blank{1.2cm}. That one fact wears several different names,
depending on who is asking.

\smallskip
\renewcommand{\arraystretch}{1.4}
\begin{center}
\begin{tabular}{|l|l|}
\hline
\textbf{The wording} & \textbf{What it asks for} \\
\hline
$a$ is a \textbf{zero} of $p$ & an input that makes the \blank{1.8cm} zero \\
$a$ is a \textbf{root} of $p(x) = 0$ & a \blank{1.8cm} of the equation \\
$(a, 0)$ is an \textbf{$x$-intercept} & where the \blank{1.6cm} meets the $x$-axis \\
\hline
\end{tabular}
\end{center}

\smallskip
\textbf{Example 1.} Is $x = 3$ a zero of $p(x) = x^3 - 2x^2 - 5x + 6$?

\begin{work}
  p(3) &= (3)^3 - 2(3)^2 - 5(3) + 6 \\
       &= 27 - 18 - 15 + 6 \\
       &= 0
\end{work}

Because $p(3) = 0$: the number $3$ is a \blank{1.6cm} of $p$, the point
\blank{1.6cm} is an $x$-intercept of the graph, and $x = 3$ is a
\blank{1.8cm} of the equation $p(x) = 0$.

\smallskip
\textit{Looking ahead:} a polynomial can only change sign at a real zero, so
these zeros become the \textit{endpoints} of the intervals we will use to
solve polynomial inequalities in Lesson 2.7.
\end{notesbox}

\vspace{0.1in}

\boxguard[30]
\begin{notesbox}{2. The Factor/Zero Equivalence}
This is the spine of the lesson. For a \textbf{real} number $a$:

\begin{center}
\fbox{\parbox{0.86\linewidth}{\centering
$(x - a)$ is a factor of $p(x)$ \quad \textbf{if and only if} \quad
$a$ is a zero of $p$.}}
\end{center}

\textbf{Reading zeros off a factored form.} Set each factor equal to zero ---
the zero-product property from the warm-up.

\smallskip
$p(x) = (x - 2)(x + 3)(x - 5)$ \qquad zeros: \blank{1cm} , \blank{1cm} ,
\blank{1cm}

\begin{itemize}[itemsep=3pt]
  \item Watch the sign flip: the factor $(x + 3)$ gives the zero
        $x = $ \blank{1.2cm} , not $+3$.
  \item This polynomial has degree \blank{1cm} and \blank{1cm} real zeros, so
        its graph has \blank{1cm} $x$-intercepts.
\end{itemize}

\textbf{Writing a factored form from zeros.} Run the equivalence backwards:
each zero $a$ contributes the factor \blank{2cm}.

\smallskip
Degree 3, zeros $-1$, $0$, and $4$: \quad $p(x) = $ \blank{5cm}

\smallskip
\textbf{Example 2.} Write a degree-2 polynomial with zeros $2$ and $-4$ whose
graph passes through $(0, -16)$.

\begin{work}
  p(x) &= a(x-2)(x+4) \\
  p(0) &= a(0-2)(0+4) \\
       &= -8a \\
   -8a &= -16 \\
     a &= 2
\end{work}

$p(x) = $ \blank{5cm}
\end{notesbox}

\vspace{0.1in}

\boxguard[34]
\begin{notesbox}{3. Multiplicity --- and What the Graph Does There}
When the linear factor $(x - a)$ is repeated $n$ times, the zero $a$ has
\textbf{multiplicity} \blank{1cm}.

\smallskip
$p(x) = (x + 1)^2(x - 3)(x - 4)^3$ \qquad degree $= 2 + 1 + 3 = $ \blank{1cm}

\smallskip
\renewcommand{\arraystretch}{1.5}
\begin{center}
\begin{tabular}{|c|c|c|c|}
\hline
\textbf{Zero} & \textbf{Multiplicity} & \textbf{Even or odd?} & \textbf{Graph at that intercept} \\
\hline
$x = -1$ & \blank{1cm} & \blank{1.2cm} & \blank{2.4cm} \\
\hline
$x = 3$  & \blank{1cm} & \blank{1.2cm} & \blank{2.4cm} \\
\hline
$x = 4$  & \blank{1cm} & \blank{1.2cm} & \blank{2.4cm} \\
\hline
\end{tabular}
\end{center}

\medskip
\textbf{The rule.}
\begin{itemize}[itemsep=4pt]
  \item \textbf{Odd} multiplicity: the outputs on the two sides of $a$ have
        \blank{2cm} signs, so the graph \blank{2cm} the $x$-axis.
  \item \textbf{Even} multiplicity: the outputs on both sides of $a$ have the
        \blank{1.6cm} sign, so the graph touches the axis and turns back ---
        it is \blank{2cm} to the $x$-axis at $x = a$.
\end{itemize}

\medskip
\textbf{Read the graph.} Below is $y = (x + 1)^2(x - 2)$.

\begin{center}
\begin{tikzpicture}[x=0.62cm, y=0.5cm]
  \draw[linegray, very thin, step=1] (-2,-2) grid (3,3);
  \draw[->, thick] (-2.3,0) -- (3.3,0) node[below right] {\small $x$};
  \draw[->, thick] (0,-2.3) -- (0,3.3) node[above] {\small $y$};
  \foreach \x in {-2,-1,1,2,3} \node[below] at (\x,-0.12) {\small $\x$};
  \draw[very thick, plum] plot [smooth] coordinates
    {(-2,-2) (-1.5,-0.44) (-1,0) (-0.5,-0.31) (0,-1) (0.5,-1.69)
     (1,-2) (1.5,-1.56) (2,0) (2.4,2.31)};
  \fill[plum] (-1,0) circle (2.5pt);
  \fill[plum] (2,0) circle (2.5pt);
\end{tikzpicture}
\end{center}

\begin{itemize}[itemsep=3pt]
  \item At $x = -1$ the graph is \blank{2cm} to the axis, matching the factor
        $(x+1)$ appearing \blank{1.4cm} times.
  \item At $x = 2$ the graph \blank{1.8cm} the axis, matching multiplicity
        \blank{1cm}.
\end{itemize}

\smallskip
\textit{Looking ahead:} here the multiplicities add up to the degree. When
they do \textit{not}, something is missing --- that is Lesson 2.4's question.
\end{notesbox}

\vspace{0.1in}

\boxguard[30]
\begin{notesbox}{4. Finding the Degree from Successive Differences}
In Lesson 2.1 a \textbf{linear} function had constant \textit{first}
differences and a \textbf{quadratic} had constant \textit{second}
differences. The pattern keeps going.

\begin{center}
\fbox{\parbox{0.90\linewidth}{\centering
Over \textbf{equally spaced} inputs, the degree of a polynomial is the
\textit{least} $n$ for which the $n$th differences are constant.}}
\end{center}

Complete each row by subtracting neighboring entries in the row above.

\smallskip
\renewcommand{\arraystretch}{1.5}
\begin{center}
\begin{tabular}{|l|c|c|c|c|c|c|}
\hline
$x$ & 0 & 1 & 2 & 3 & 4 & 5 \\
\hline
$f(x)$ & 5 & 4 & 9 & 26 & 61 & 120 \\
\hline
1st differences & \multicolumn{6}{c|}{\blank{1.2cm} \ \blank{1.2cm} \ \blank{1.2cm} \ \blank{1.2cm} \ \blank{1.2cm}} \\
\hline
2nd differences & \multicolumn{6}{c|}{\blank{1.2cm} \ \blank{1.2cm} \ \blank{1.2cm} \ \blank{1.2cm}} \\
\hline
3rd differences & \multicolumn{6}{c|}{\blank{1.2cm} \ \blank{1.2cm} \ \blank{1.2cm}} \\
\hline
\end{tabular}
\end{center}

\begin{itemize}[itemsep=3pt]
  \item The first constant row is the \blank{1.4cm} differences, so $f$ has
        degree \blank{1cm}.
  \item Careful: this test only works when the inputs are \blank{3.4cm}.
\end{itemize}
\end{notesbox}

\vspace{0.1in}

\boxguard
\begin{practicebox}
\begin{enumerate}[itemsep=8pt]
  \item Let $p(x) = (x - 6)(x + 2)^2$. Its degree is \blank{1cm}.

        \begin{enumerate}[label=(\alph*), itemsep=4pt]
          \item Zero $x = 6$, multiplicity \blank{1cm}; the graph
                \blank{2.4cm} the $x$-axis there.
          \item Zero $x = $ \blank{1.2cm} , multiplicity \blank{1cm}; the
                graph is \blank{2.4cm} to the $x$-axis there.
        \end{enumerate}

  \item Write a degree-3 polynomial in factored form with a zero at $x = 3$ of
        multiplicity 2, a zero at $x = -5$, and $p(0) = 90$.

        \begin{work}
          p(x) &= a(x-3)^2(x+5) \\
          p(0) &= a(-3)^2(5) \\
               &= 45a \\
           45a &= 90 \\
             a &= 2
        \end{work}

        $p(x) = $ \blank{5cm}
\end{enumerate}
\end{practicebox}

\end{document}
